% !TEX root = ../../main.tex
% C3.2 — Phát hiện người trong ảnh (RÚT GỌN 2026-09-29; bản cũ ở report/archive/)
% Nguyên lý; lý do chọn YOLO11n so với phương án khác để ở Mục 6.1.4.
% Khớp app: config/ultralytics_yolo_detector.json (conf 0.1 = track_low_threshold của ByteTrack, iou 0.7, max 300), chỉ lớp person, confidence chỉ dùng nội bộ.
\section{Phát hiện người trong ảnh}
\label{sec:3.2}

\subsection{Bài toán phát hiện đối tượng}
\label{sec:3.2.1}

Phát hiện đối tượng (object detection) xác định vị trí và loại của các đối tượng trong một ảnh. Với mỗi đối tượng, mô hình trả về một \emph{khung bao} (bounding box) hình chữ nhật, thường biểu diễn bằng tọa độ hai góc $(x_1, y_1, x_2, y_2)$, \emph{lớp} của đối tượng và một \emph{độ tin cậy} (confidence score). Hai khung bao được so sánh bằng chỉ số giao trên hợp (Intersection over Union, IoU):
\begin{equation}
    \mathrm{IoU}(A, B) = \frac{|A \cap B|}{|A \cup B|},
    \label{eq:iou}
\end{equation}
nhận giá trị từ 0 khi hai khung không giao nhau đến 1 khi trùng khít. IoU được dùng khi loại bỏ các dự đoán trùng lặp và khi liên kết phát hiện giữa các khung hình trong bài toán theo vết (Mục~\ref{sec:3.3}).

\subsection{Hướng tiếp cận hai giai đoạn và một giai đoạn}
\label{sec:3.2.2}

Mô hình \emph{hai giai đoạn} như Faster R-CNN~\cite{ren2017fasterrcnn} trước tiên dùng một mạng đề xuất vùng để sinh các vùng có khả năng chứa đối tượng, sau đó phân loại từng vùng và tinh chỉnh khung bao; cách này chính xác nhưng tốn kém do phải xử lý nhiều vùng đề xuất. Mô hình \emph{một giai đoạn} như YOLO (You Only Look Once)~\cite{redmon2016yolo} xem phát hiện là một bài toán hồi quy duy nhất, dự đoán trực tiếp khung bao và lớp của mọi đối tượng trong một lần lan truyền qua ảnh, nên nhanh hơn đáng kể và phù hợp với việc xử lý nhiều khung hình video.

\subsection{Họ mô hình YOLO}
\label{sec:3.2.3}

Các mô hình YOLO hiện đại gồm ba phần: phần trích xuất đặc trưng (backbone) ở nhiều mức chi tiết, phần kết hợp đặc trưng (neck) ở nhiều tỉ lệ để phát hiện được cả đối tượng lớn gần camera lẫn đối tượng nhỏ ở xa, và phần dự đoán (head) cho khung bao, lớp và độ tin cậy. YOLO11 do Ultralytics phát hành năm 2024 với năm kích cỡ (n, s, m, l, x) để đánh đổi giữa tốc độ và độ chính xác~\cite{ultralytics2024yolo11docs}. Các mô hình được huấn luyện sẵn trên COCO~\cite{lin2014coco}, gồm 80 lớp đối tượng thông dụng trong đó có lớp \emph{người}, nên dùng được trực tiếp để phát hiện người mà không cần huấn luyện thêm.

\subsection{Hậu xử lý: ngưỡng tin cậy và loại bỏ trùng lặp}
\label{sec:3.2.4}

Đầu ra thô có nhiều khung độ tin cậy thấp và nhiều khung chồng lấn. Khung dưới ngưỡng tin cậy bị loại, và phép loại bỏ khung không cực đại (Non-Maximum Suppression, NMS) giữ khung tin cậy nhất, bỏ các khung cùng lớp chồng lấn với nó quá ngưỡng IoU. Ngưỡng tin cậy thấp ít bỏ sót nhưng thêm phát hiện sai; ngưỡng IoU thấp loại chồng lấn mạnh hơn nhưng có thể gộp nhầm người đứng sát nhau.

\subsection{Vai trò của bước phát hiện trong hệ thống}
\label{sec:3.2.5}

Mô hình phát hiện được áp dụng lên các khung hình đã được lấy mẫu (Mục~\ref{sec:3.4}) và chỉ giữ lớp \emph{người}, với ngưỡng tin cậy 0,1, ngưỡng IoU của NMS 0,7 và tối đa 300 phát hiện trên một khung hình. Ngưỡng thấp là chủ ý: ByteTrack cần các phát hiện độ tin cậy thấp, thường là người bị che khuất một phần, để giữ track (Mục~\ref{sec:3.3}); chúng không khởi tạo track, và track cần hai lần quan sát mới được xác nhận, nên phần lớn phát hiện sai thoáng qua bị loại ở bước theo vết. Lý do chọn YOLO11n được trình bày tại Mục~\ref{sec:6.1.4}.
